\begin{table}[h!]
    \begin{tabular}{| p{3cm} | p{9cm} |}
    \hline
    \textbf{Kode} 
    & \textbf{Deskripsi Kebutuhan} \\
    \hline

    \textbf{NFR1-Reproducibility}
    & \textit{pipeline} dirancang sehingga eksperimen dapat diulang dengan hasil yang sama melalui pengaturan random seed, dokumentasi konfigurasi, dan struktur eksekusi yang deterministik. \\
    \hline
    
    \textbf{NFR2-Efisiensi Komputasi} 
    & \textit{pipeline} dapat dijalankan pada perangkat komputasi standar (CPU tanpa GPU) dengan waktu pemrosesan yang wajar, sehingga prosedur komparasi enam model tidak membebani sumber daya.\\

    \hline
    \textbf{NFR3-Skalabilitas Moderat} 
    & \textit{pipeline} harus tetap stabil ketika jumlah data meningkat (misalnya dari 500 menjadi >1000 artikel), tanpa memerlukan perubahan arsitektural signifikan atau penurunan performa yang drastis. \\

    \hline
    \textbf{NFR4-Kemudahan Konfigurasi} 
    & \textit{pipeline} harus memungkinkan pengaturan parameter dari model (misalnya jumlah pohon pada Random Forest) dan parameter fitur (misalnya n-gram batas TF-IDF) melalui modul konfigurasi yang mudah diubah. \\

    \hline
    \textbf{NFR5-Kemudahan Interpretasi Hasil} 
    & \textit{pipeline} harus menghasilkan laporan evaluasi yang ringkas, terstruktur, dan mudah dipahami, sehingga hasil eksperimen dapat dijelaskan kepada pembimbing atau pihak lain dengan penjelasan tersebut. \\
    \hline
    \end{tabular}
\caption{Kebutuhan Non-Fungsional dalam Implementasi \textit{pipeline}}
\label{tbl:kebutuhan-non-fungsional}
\end{table}